\scanpage{149}
\mspage{23}
% archive sheet number: 148

Venons-en à la catégorie $M^r$ elle $=$. L'argument
de la page 21 montre que l'ens des \inserted{classes d'isom.}
objets de $M^r$ indexés par un $I$ fixé, de cardinal $\leq r^{\pi_0}$, comme
$\deleted{\unsure{X}}A=\underset{i}{\underline{\lim}}\,A_i$ avec $A_i\in\unsure{\mathrm{Filt}}^{\pi_0}$, est de cardinal
majoré par $r^{\pi_0}$ (si $r\geq C_0$). Si on fait varier
$I$, on trouve \deleted{\unsure{$r^{\pi}$}} une majoration par $r^{\pi_0}\mathbin{\cdot}2^{(r^{\pi_0})}
=2^{(r^{\pi_0})}$. On en conclut comme tantôt que
l'ens $\widetilde{\mathrm{Fl}\, M^r}$ est de cardinal majoré par
\[
\left(2^{(r^{\pi_0})}\right)^2\times r^{(r^{\pi_0})}
\qquad(\text{si }r\geq\deleted{\unsure{\pi}}\operatorname{Sup}(\pi_1,C_1)),
\]
donc par $r^{(r^{\pi_0})}$. Je dis qu'on a
\[
r^{(r^{\pi_0})}=2^{(r^{\pi_0})},
\]
en effet, on a $\geq$, et pour prouver $\leq$, on
note que $r<2^r$, donc $r^{(r^{\pi_0})}\leq(2^r)^{(r^{\pi_0})}=2^{r\mathbin{\cdot}r^{\pi_0}}=2^{r^{\pi_0}}$,
OK ! Cela établit donc 3.7. \deleted{On en déduit} \inserted{Cet argument}
\unsure{donne} en fait le

\uline{Lemme 3.\unsure{?|2}.} Soient $r$ un cardinal, $r'$ un cardinal $\geq r$. Alors
on a
\[
r^{r'}=2^{r'}\qquad(=s^{r'}\text{ pour tout }2\leq s\leq r)
\]

\uline{Corollaire 1} Soient \deleted{$r\geq s\geq\operatorname{Sup}(\pi_1,C_1)$.} $r\geq\operatorname{Sup}(\pi_1,C_1)$,
et $A,X\in\unsure{\mathrm{Pfilt}|\mathrm{Filt}}^{\unsure{\pi|r}}$. Alors
\[
\operatorname{Card}\operatorname{Hom}_M(A,X)\leq 2^{(r^{\pi_0})}
\qquad(=2^r\text{ si }r=2^c,\ \inserted{\text{avec }}c\geq\pi_0)
\tag*{(3.10.),}
\]
On applique 3.8., \deleted{\unsure{avec}} qui donne la majoration
$r^{(r^{\pi_0})}$, et \unsure{le|les} lemme 3.2.

\deleted{\uline{Corollaire 2}} On notera qu'on trouve la même
majoration que pour le cardinal de $\widetilde{\underline{\mathrm{Fl}}(M^r)}$
\unsure{dont mention} (3.7) ! De $=$ (3.9.) donne

\uline{Corollaire 2} Soit $r\geq\operatorname{Sup}(\pi_1,C_1)$, et $\unsure{X|\alpha}\in\operatorname{Ob}M^r$. Alors
\[
\operatorname{Card}\widetilde{\underline{\mathrm{Fl}}(M^r_{/X})}\leq 2^{(\unsure{r}^{\pi_0})}
\qquad(=2^r\text{ si }r=2^c,\text{ avec }c\geq\pi_0).
\tag*{(3.11).}
\]
c'est donc la formule (3.9) pour $s\leq r$, lorsque on
suppose.
